\chapter{Krylov Spaces and Random Sketches}\label{ch:krylov}

\begin{center}\itshape
The method of conjugate gradients was developed independently by E.~Stiefel of the Institute of Applied Mathematics\\
at Z\"urich and by M.~R.~Hestenes with the cooperation of J.~B.~Rosser, G.~Forsythe, and L.~Paige.\\
\upshape(Hestenes and Stiefel, 1952)
\end{center}

In the summer of 1951 the Institute for Numerical Analysis, a branch of the National Bureau of Standards housed on the campus of the University of
California in Los Angeles, was one of the few places in the world with an electronic computer, the SWAC, and a crowd of mathematicians trying to
find out what to do with it. One of them, Magnus Hestenes of UCLA, had come to linear systems from the calculus of variations: solving $Ax=b$ for a
symmetric positive definite $A$ is minimising the quadratic $\frac12x^{\mathsf T}Ax-b^{\mathsf T}x$, and minimising is what variational people do.
Eduard Stiefel, visiting from the ETH in Z\"urich, had arrived at the same algorithm from a quite different direction. At a symposium in August 1951
Stiefel and Barkley Rosser reported on it, and the joint paper of Hestenes and Stiefel, \emph{Methods of conjugate gradients for solving linear
systems}, appeared in December 1952. Cornelius Lanczos, at the same institute, had found a closely related method for eigenvalues.

The method had a curious fate. In exact arithmetic it finds the solution of an $n\times n$ system in at most $n$ steps, so it was first seen as a
direct method competing with elimination, and in that role it lost: rounding errors spoil the finite termination, and naive implementations on the
large problems of the 1960s were disappointing. It went into eclipse. Then in 1971 John Reid observed that on large sparse systems the method gave
excellent answers \emph{long before} step $n$, and it came back as an iterative method, which is what it really is. In 2000 the journal \emph{Computing in
Science and Engineering} listed Krylov subspace iteration among the ten algorithms of the century, next to the $QR$ algorithm of Chapter~\ref{ch:matrices}.

This chapter explains why it converges long before step $n$: the $k$th iterate is the best approximation in a space of polynomials in $A$, and the
best polynomial adapts itself to the spectrum. The same Krylov spaces give the Lanczos method for eigenvalues, which we use for spectral clustering,
keeping a promise of Tome~C1, Chapter~9. The second half of the chapter makes matrices small by a different trick, randomness: a Gaussian random
matrix with a few more rows than a subspace has dimensions preserves all lengths in it, and this makes least squares and low-rank approximation of
huge matrices cheap. The tools are $\varepsilon$-nets and the chi-square tail bounds of Tome~10, Chapter~8, where the Johnson--Lindenstrauss lemma was
proved; here we upgrade it from finitely many points to whole subspaces.

\begin{think}
The matrix of the discrete Laplacian on a $100\times100$ grid is $10^4\times10^4$, symmetric positive definite, with condition number about
$4\times10^3$. Gaussian elimination ignores its sparsity and costs about $10^{12}$ operations. A method that needs only products $Ax$, each costing
$5\times10^4$ operations, would win if it needed fewer than $10^7$ of them. How many do you think conjugate gradients needs to gain six digits?
\end{think}

\section{Krylov spaces and Chebyshev polynomials}\label{sec:krylovspaces}

\begin{definition}\label{def:krylov}
For an $n\times n$ matrix $A$, a vector $b$ and $k\ge1$, the \emph{Krylov space} is $\mathcal K_k(A,b)=\operatorname{span}(b,Ab,\dots,A^{k-1}b)=\{q(A)b:q\in\mathcal P_{k-1}\}$,
where $\mathcal P_j$ is the space of real polynomials of degree at most $j$.
\end{definition}

Krylov spaces are the natural search spaces of any method that only multiplies by $A$: after $k-1$ products, all we know lies in $\mathcal K_k$. The
quality of the best element of $\mathcal K_k$ is measured by a polynomial approximation problem, and for problems with eigenvalues in an interval
the answer to such problems is given by Chebyshev polynomials. Recall that $T_0=1$, $T_1(x)=x$, $T_{j+1}(x)=2xT_j(x)-T_{j-1}(x)$, so $T_j$ has degree $j$.

\begin{lemma}[Chebyshev polynomials]\label{lem:chebyshev}
For $x\in[-1,1]$, $T_j(x)=\cos(j\arccos x)$, so $|T_j(x)|\le1$. For $x\ge1$,
\[
T_j(x)=\tfrac12\Bigl[\bigl(x+\sqrt{x^2-1}\bigr)^j+\bigl(x-\sqrt{x^2-1}\bigr)^j\Bigr]\ge\tfrac12\bigl(x+\sqrt{x^2-1}\bigr)^j.
\]
\end{lemma}
\begin{proof}
Both right-hand sides satisfy the recurrence and the initial values: for the first, $\cos((j+1)t)+\cos((j-1)t)=2\cos t\cos(jt)$; for the second, $z_\pm=x\pm\sqrt{x^2-1}$
are the roots of $z^2=2xz-1$, so $z_\pm^{j+1}=2xz_\pm^j-z_\pm^{j-1}$. The last inequality holds because $z_->0$ when $x\ge1$.
\end{proof}

If the eigenvalues of a symmetric positive definite $A$ lie in $[\mu,L]$ with $0<\mu<L$, put $\kappa=L/\mu$. The affine map $\lambda\mapsto(L+\mu-2\lambda)/(L-\mu)$
sends $[\mu,L]$ onto $[-1,1]$ and $0$ to $x_0=\frac{\kappa+1}{\kappa-1}$, for which $x_0+\sqrt{x_0^2-1}=\frac{\sqrt\kappa+1}{\sqrt\kappa-1}$ (square it). Hence the polynomial
\begin{equation}\label{eq:chebpoly}
p_k(\lambda)=\frac{T_k\bigl((L+\mu-2\lambda)/(L-\mu)\bigr)}{T_k(x_0)}
\end{equation}
satisfies $p_k(0)=1$ and
\[
\max_{\lambda\in[\mu,L]}|p_k(\lambda)|\le\frac1{T_k(x_0)}\le2\Bigl(\frac{\sqrt\kappa-1}{\sqrt\kappa+1}\Bigr)^k.
\]
Exercise~\ref{ex:chebopt} shows that no polynomial of degree $k$ with value $1$ at $0$ does better on the whole interval. The quantity
$\frac{\sqrt\kappa-1}{\sqrt\kappa+1}\approx1-\frac2{\sqrt\kappa}$ is the hallmark of the methods of this chapter and of Nesterov's method in
Chapter~4; the cruder $\frac{\kappa-1}{\kappa+1}\approx1-\frac2\kappa$ belongs to methods that do not exploit the polynomial freedom (Tome~S1, Theorem~12.6 compares the
two for stationary iterations).

\section{Conjugate gradients}\label{sec:cg}

Let $A$ be symmetric positive definite, $x_*=A^{-1}b$, and write $\|y\|_A=(y^{\mathsf T}Ay)^{1/2}$ for the \emph{energy norm}. Since
$\frac12x^{\mathsf T}Ax-b^{\mathsf T}x=\frac12\|x-x_*\|_A^2-\frac12x_*^{\mathsf T}Ax_*$, minimising the quadratic over a set is minimising the energy distance to $x_*$. We
start from $x_0=0$ (otherwise replace $b$ by $b-Ax_0$).

\begin{definition}\label{def:cg}
The \emph{conjugate gradient iterate} $x_k$ is the minimiser of $\|x-x_*\|_A$ over $x\in\mathcal K_k(A,b)$.
\end{definition}

Since $\mathcal K_k$ is a subspace and $\|\cdot\|_A$ a Euclidean norm, $x_k$ exists and is characterised by the orthogonality $(x_k-x_*)^{\mathsf T}Av=0$ for all
$v\in\mathcal K_k$, that is, the \emph{residual} $r_k=b-Ax_k$ is orthogonal to $\mathcal K_k$ (projection theorem, Tome~S1, Theorem~4.7). This is a definition,
not an algorithm; the miracle is that $x_k$ can be updated with a constant amount of work and storage per step.

\begin{theorem}[Hestenes and Stiefel, 1952]\label{thm:cg}
Put $r_0=p_0=b$, and for $k=1,2,\dots$, as long as $r_{k-1}\ne0$,
\[
\begin{gathered}
\alpha_k=\frac{r_{k-1}^{\mathsf T}r_{k-1}}{p_{k-1}^{\mathsf T}Ap_{k-1}},\qquad\beta_k=\frac{r_k^{\mathsf T}r_k}{r_{k-1}^{\mathsf T}r_{k-1}},\\
x_k=x_{k-1}+\alpha_kp_{k-1},\qquad r_k=r_{k-1}-\alpha_kAp_{k-1},\qquad p_k=r_k+\beta_kp_{k-1}.
\end{gathered}
\]
Then $r_k=b-Ax_k$, the vectors $x_k$ are the conjugate gradient iterates of Definition~\ref{def:cg}, the residuals are orthogonal, $r_j^{\mathsf T}r_i=0$ for $i\ne j$, the
\emph{search directions} are $A$-conjugate, $p_j^{\mathsf T}Ap_i=0$ for $i\ne j$, and
$\mathcal K_k(A,b)=\operatorname{span}(r_0,\dots,r_{k-1})=\operatorname{span}(p_0,\dots,p_{k-1})$.
\end{theorem}
\begin{proof}
$r_k=b-Ax_k$ follows from the updates by induction. We prove by induction on $k$ that, as long as $r_0,\dots,r_{k-1}\ne0$: (a) $r_0,\dots,r_{k-1}$ and
$p_0,\dots,p_{k-1}$ span $\mathcal K_k$; (b) $r_j^{\mathsf T}r_i=0$ for $i<j\le k$; (c) $p_j^{\mathsf T}Ap_i=0$ for $i<j\le k$. For $k=1$, (a) is clear.

\emph{Spans.} By the updates, $r_k$ is a combination of $r_{k-1}$ and $Ap_{k-1}$, and $p_k$ a combination of $r_k$ and $p_{k-1}$; so both lie in $\mathcal K_{k+1}$. Once (b) holds, the
nonzero vectors $r_0,\dots,r_{k-1}$ are orthogonal, hence independent, and $k$ independent vectors in $\mathcal K_k$, of dimension at most $k$, span it; the same
holds for the $p_i$, since $p_i-r_i\in\operatorname{span}(p_{i-1})$.

\emph{Residuals.} Let $i<j$. Then $r_j^{\mathsf T}r_i=r_{j-1}^{\mathsf T}r_i-\alpha_jp_{j-1}^{\mathsf T}Ar_i$, and $r_i=p_i-\beta_ip_{i-1}$ (with $p_{-1}=0$). If $i<j-1$, both terms
vanish by induction: $r_{j-1}\perp r_i$, and $p_{j-1}^{\mathsf T}Ar_i=p_{j-1}^{\mathsf T}Ap_i-\beta_ip_{j-1}^{\mathsf T}Ap_{i-1}=0$. If $i=j-1$,
$p_{j-1}^{\mathsf T}Ar_{j-1}=p_{j-1}^{\mathsf T}Ap_{j-1}$ by conjugacy, and $r_j^{\mathsf T}r_{j-1}=r_{j-1}^{\mathsf T}r_{j-1}-\alpha_jp_{j-1}^{\mathsf T}Ap_{j-1}=0$ by the choice of $\alpha_j$.

\emph{Directions.} Let $i<j$. Then $p_j^{\mathsf T}Ap_i=r_j^{\mathsf T}Ap_i+\beta_jp_{j-1}^{\mathsf T}Ap_i$, and $Ap_i=(r_i-r_{i+1})/\alpha_{i+1}$. If $i<j-1$, the second term vanishes by
induction and the first is $r_j^{\mathsf T}(r_i-r_{i+1})/\alpha_{i+1}=0$ by (b). If $i=j-1$, $r_j^{\mathsf T}Ap_{j-1}=-r_j^{\mathsf T}r_j/\alpha_j$ and
$p_{j-1}^{\mathsf T}Ap_{j-1}=r_{j-1}^{\mathsf T}r_{j-1}/\alpha_j$, so $p_j^{\mathsf T}Ap_{j-1}=(-r_j^{\mathsf T}r_j+\beta_jr_{j-1}^{\mathsf T}r_{j-1})/\alpha_j=0$ by the choice of $\beta_j$.

The divisions are legitimate: $p_{k-1}^{\mathsf T}Ap_{k-1}>0$ since $p_{k-1}\ne0$, because $p_{k-1}^{\mathsf T}r_{k-1}=r_{k-1}^{\mathsf T}r_{k-1}>0$ (as $p_{k-1}-r_{k-1}\in\mathcal K_{k-1}\perp r_{k-1}$).
Finally $x_k\in\operatorname{span}(p_0,\dots,p_{k-1})=\mathcal K_k$ and $r_k\perp r_0,\dots,r_{k-1}$, that is $r_k\perp\mathcal K_k$, which characterises the conjugate gradient iterate.
\end{proof}

Each step costs one product $Ap_{k-1}$, two inner products and three vector updates, and only four vectors are stored: for a sparse $A$ with $s$ nonzero
entries, $O(s+n)$ per step. Now the convergence. With $x_0=0$, every $x\in\mathcal K_k$ is $q(A)b=q(A)Ax_*$ with $q\in\mathcal P_{k-1}$, so
$x-x_*=-(I-Aq(A))x_*=p(A)(x_0-x_*)$ with $p(\lambda)=1-\lambda q(\lambda)$: the errors of the Krylov space are the vectors $p(A)e_0$, $p\in\mathcal P_k$, $p(0)=1$.

\begin{theorem}[convergence of conjugate gradients]\label{thm:cgconv}
Let $\lambda_1,\dots,\lambda_n$ be the eigenvalues of $A$, and $e_k=x_k-x_*$. Then
\[
\|e_k\|_A=\min_{p\in\mathcal P_k,\,p(0)=1}\|p(A)e_0\|_A\le\min_{p\in\mathcal P_k,\,p(0)=1}\ \max_i|p(\lambda_i)|\cdot\|e_0\|_A.
\]
Consequently:
\begin{enumerate}
\item if $A$ has $m$ distinct eigenvalues, $x_m=x_*$;
\item if the eigenvalues lie in $[\mu,L]$, $\kappa=L/\mu$, then $\|e_k\|_A\le2\bigl(\frac{\sqrt\kappa-1}{\sqrt\kappa+1}\bigr)^k\|e_0\|_A$;
\item if all eigenvalues but $\lambda_{i_1},\dots,\lambda_{i_m}$ lie in $[\mu,L]$ and the exceptional ones exceed $L$, then for $k\ge m$,
$\|e_k\|_A\le2\bigl(\frac{\sqrt{\kappa}-1}{\sqrt{\kappa}+1}\bigr)^{k-m}\|e_0\|_A$ with $\kappa=L/\mu$.
\end{enumerate}
\end{theorem}
\begin{proof}
The equality is Definition~\ref{def:cg} rewritten. Expand $e_0=\sum_ic_iv_i$ in an orthonormal eigenbasis; then $\|p(A)e_0\|_A^2=\sum_ic_i^2\lambda_ip(\lambda_i)^2\le\max_ip(\lambda_i)^2\sum_ic_i^2\lambda_i$,
which is the inequality. (1) Take $p(\lambda)=\prod_{j=1}^m(1-\lambda/\theta_j)$ over the distinct eigenvalues $\theta_j$. (2) Take $p_k$ of~\eqref{eq:chebpoly}. (3) Take
$p(\lambda)=p_{k-m}(\lambda)\prod_{l=1}^m(1-\lambda/\lambda_{i_l})$: it vanishes at the exceptional eigenvalues, and on $[\mu,L]$ each factor $1-\lambda/\lambda_{i_l}$ lies in $[0,1]$.
\end{proof}

Part (3) is the promise of Tome~S1, Chapter~12: the method was never told about the outliers, yet after $m$ steps it behaves as if they had been removed,
because the optimal polynomial places roots on them by itself. Clustered spectra give the same effect, and this is the \emph{superlinear} convergence
seen in practice. For the Laplacian of the \emph{Pause and think}, $\kappa\approx4134$, $\sqrt\kappa\approx64$, and (2) promises six digits in energy norm after about
$\frac{\sqrt\kappa}2\ln(2\times10^6)\approx467$ steps; the notebook needs $264$ (Figure~\ref{fig:cgconv}). Gaussian elimination would have
cost $10^{12}$ operations, conjugate gradients about $4\times10^7$.

\begin{figure}[ht]
\centering
\fig{ch02-cgconv}
\caption{Conjugate gradients on the discrete Laplacian of a $100\times100$ grid ($n=10^4$, $\kappa\approx4.1\times10^3$, right-hand side with random entries): the relative error in energy norm against the
bound $2\bigl(\frac{\sqrt\kappa-1}{\sqrt\kappa+1}\bigr)^k$ of Theorem~\ref{thm:cgconv}, and steepest descent, whose worst-case rate is $\frac{\kappa-1}{\kappa+1}$ per step.}
\label{fig:cgconv}
\end{figure}

\begin{remark}
Two practical facts we do not prove. In floating point the residuals soon lose their orthogonality and finite termination fails, but convergence
continues (Anne Greenbaum, 1989, showed that finite precision CG behaves like exact CG on a larger matrix whose eigenvalues are clustered in tiny intervals
around those of $A$). Usually the predicted rate survives; a huge isolated eigenvalue, however, then counts several times, and the notebook shows five outliers
of size up to $10^6$ costing four hundred extra steps instead of fifteen. And everything depends on $\kappa$, so one solves $M^{-1}Ax=M^{-1}b$ instead, with a \emph{preconditioner} $M$
that approximates $A$ and is easy to invert; finding good preconditioners is an art, and for graph Laplacians it became a theorem: Spielman and Teng
proved in 2004 that symmetric diagonally dominant systems can be solved in nearly linear time.
\end{remark}

\section{Lanczos and the extreme eigenvalues}\label{sec:lanczos}

Orthonormalising $b,Ab,A^2b,\dots$ by Gram--Schmidt gives an orthonormal basis $q_1,q_2,\dots$ of the Krylov spaces. For symmetric $A$, the process is
short.

\begin{proposition}[Lanczos, 1950]\label{prop:lanczos}
Let $A$ be symmetric, $q_0=0$, $q_1=b/\|b\|$, and for $j\ge1$: $\alpha_j=q_j^{\mathsf T}Aq_j$, $w_j=Aq_j-\alpha_jq_j-\beta_{j-1}q_{j-1}$, $\beta_j=\|w_j\|$,
$q_{j+1}=w_j/\beta_j$ (with $\beta_0=0$), as long as $\beta_j\ne0$. Then $q_1,\dots,q_k$ is an orthonormal basis of $\mathcal K_k(A,b)$, and with $Q_k=(q_1\cdots q_k)$,
\[
AQ_k=Q_kT_k+\beta_kq_{k+1}e_k^{\mathsf T},\qquad T_k=Q_k^{\mathsf T}AQ_k=\begin{pmatrix}\alpha_1&\beta_1\\\beta_1&\alpha_2&\ddots\\&\ddots&\ddots&\beta_{k-1}\\&&\beta_{k-1}&\alpha_k\end{pmatrix}.
\]
\end{proposition}
\begin{proof}
Gram--Schmidt on $q_1,Aq_1,Aq_2,\dots$ produces an orthonormal basis of the Krylov spaces, because $Aq_j\in\mathcal K_{j+1}$ and, inductively, $Aq_j\notin\mathcal K_j$ while
$\beta_j\ne0$. The component of $Aq_j$ along $q_i$ is $q_i^{\mathsf T}Aq_j=(Aq_i)^{\mathsf T}q_j$, which vanishes for $i<j-1$ because $Aq_i\in\mathcal K_{i+1}\perp q_j$. So Gram--Schmidt only
needs to remove the components along $q_j$ and $q_{j-1}$, the latter being $q_{j-1}^{\mathsf T}Aq_j=(Aq_{j-1})^{\mathsf T}q_j=\beta_{j-1}$ (as
$Aq_{j-1}=\beta_{j-1}q_j+\alpha_{j-1}q_{j-1}+\beta_{j-2}q_{j-2}$). This is the stated recurrence, and it reads $AQ_k=Q_kT_k+\beta_kq_{k+1}e_k^{\mathsf T}$ column by column;
multiplying by $Q_k^{\mathsf T}$ gives $T_k=Q_k^{\mathsf T}AQ_k$.
\end{proof}

The eigenvalues $\theta_1\ge\dots\ge\theta_k$ of $T_k$ are the \emph{Ritz values}. Since $T_k$ is $A$ compressed to $\mathcal K_k$, the Courant--Fischer
theorem (Tome~S1, Chapter~5) gives $\theta_1=\max_{0\ne x\in\mathcal K_k}R(x)$ with $R(x)=x^{\mathsf T}Ax/x^{\mathsf T}x$, so $\theta_1\le\lambda_1$; and, the
Krylov space containing the power iterates, $\theta_1$ is at least as good as the Rayleigh quotient of the power method. It is much better.

\begin{theorem}[Kaniel, 1966; Paige, 1971; Saad, 1980]\label{thm:kaniel}
Let $A$ be symmetric with eigenvalues $\lambda_1>\lambda_2\ge\dots\ge\lambda_n$ and orthonormal eigenvectors $v_i$, and let $\phi$ be the angle between $b$ and $v_1$. Then
\[
0\le\lambda_1-\theta_1\le(\lambda_1-\lambda_n)\frac{\tan^2\phi}{T_{k-1}(1+2\gamma)^2},\qquad\gamma=\frac{\lambda_1-\lambda_2}{\lambda_2-\lambda_n}.
\]
\end{theorem}
\begin{proof}
For any $p\in\mathcal P_{k-1}$ with $p(\lambda_1)\ne0$, $x=p(A)b\in\mathcal K_k$. Write $b=\sum_ic_iv_i$; then
\[
\lambda_1-R(x)=\frac{\sum_{i\ge2}c_i^2p(\lambda_i)^2(\lambda_1-\lambda_i)}{\sum_ic_i^2p(\lambda_i)^2}\le(\lambda_1-\lambda_n)\frac{\sum_{i\ge2}c_i^2}{c_1^2}\cdot\max_{i\ge2}\frac{p(\lambda_i)^2}{p(\lambda_1)^2},
\]
and $\sum_{i\ge2}c_i^2/c_1^2=\tan^2\phi$. Take $p(\lambda)=T_{k-1}\bigl(1+2\frac{\lambda-\lambda_2}{\lambda_2-\lambda_n}\bigr)$: the argument runs over $[-1,1]$ for $\lambda\in[\lambda_n,\lambda_2]$,
so $|p(\lambda_i)|\le1$ for $i\ge2$ (Lemma~\ref{lem:chebyshev}), while $p(\lambda_1)=T_{k-1}(1+2\gamma)$. Since $\theta_1\ge R(x)$, the bound follows.
\end{proof}

By Lemma~\ref{lem:chebyshev}, for small $\gamma$,
\[
T_{k-1}(1+2\gamma)\ge\tfrac12\bigl(1+2\gamma+2\sqrt{\gamma+\gamma^2}\bigr)^{k-1}\ge\tfrac12e^{2(k-1)\sqrt\gamma}:
\]
the error decays like $e^{-4k\sqrt\gamma}$,
against $(1+\gamma)^{-2k}\approx e^{-2k\gamma}$ for the power method. A relative gap of $\gamma=10^{-4}$ costs the power method about $50\,000$ steps per factor $e^{-10}$ and Lanczos
at most about $250$; the notebook measures $49\,000$ and $110$. The same holds at the other end of the spectrum, so Lanczos finds the extreme eigenvalues of huge sparse matrices with a few hundred products.
In floating point the $q_j$ lose orthogonality as soon as a Ritz value converges, and spurious copies (``ghosts'') of converged eigenvalues appear;
Christopher Paige showed in 1971 that this loss is precisely tied to convergence, and the cure is selective reorthogonalisation (Exercise~\ref{ex:ghosts}).

\section{Spectral clustering}\label{sec:spectralclustering}

Here is a use for the second smallest eigenvector of a large sparse matrix. A data set is turned into a weighted graph: the points are the vertices,
and $w_{ij}=w_{ji}\ge0$ measures the similarity of points $i$ and $j$, for example $w_{ij}=e^{-\|x_i-x_j\|^2/2s^2}$ for the $10$ nearest neighbours and
$0$ otherwise. Let $d_i=\sum_jw_{ij}>0$ be the degrees, $D=\operatorname{diag}(d_i)$, $L=D-W$ the Laplacian (so $x^{\mathsf T}Lx=\sum_{i<j}w_{ij}(x_i-x_j)^2$, as in Tome~C1, Chapter~9),
and for $S\subseteq V$ put $\operatorname{vol}S=\sum_{i\in S}d_i$ and $\operatorname{cut}(S,\bar S)=\sum_{i\in S,j\notin S}w_{ij}$. In 2000, to segment images into objects,
Jianbo Shi and Jitendra Malik proposed to minimise the \emph{normalised cut}
\[
\operatorname{Ncut}(S)=\operatorname{cut}(S,\bar S)\Bigl(\frac1{\operatorname{vol}S}+\frac1{\operatorname{vol}\bar S}\Bigr),
\]
which penalises cuts that are cheap only because one side is tiny. Minimising it exactly is NP-hard; but it is a Rayleigh quotient in disguise.

\begin{theorem}[Shi and Malik, 2000]\label{thm:ncut}
For $\emptyset\ne S\ne V$, the vector $x_S=\mathbf 1_S/\operatorname{vol}S-\mathbf 1_{\bar S}/\operatorname{vol}\bar S$ satisfies $x_S^{\mathsf T}D\mathbf1=0$ and
$\operatorname{Ncut}(S)=x_S^{\mathsf T}Lx_S/x_S^{\mathsf T}Dx_S$. Consequently, if $\lambda_2$ is the second smallest eigenvalue of the generalised problem $Lx=\lambda Dx$,
\[
\lambda_2=\min_{x\ne0,\ x^{\mathsf T}D\mathbf1=0}\frac{x^{\mathsf T}Lx}{x^{\mathsf T}Dx}\le\min_S\operatorname{Ncut}(S).
\]
\end{theorem}
\begin{proof}
$x_S^{\mathsf T}D\mathbf1=\operatorname{vol}S/\operatorname{vol}S-\operatorname{vol}\bar S/\operatorname{vol}\bar S=0$. In $x_S^{\mathsf T}Lx_S=\sum_{i<j}w_{ij}(x_i-x_j)^2$ only the pairs across the cut contribute, each
with $(x_i-x_j)^2=(1/\operatorname{vol}S+1/\operatorname{vol}\bar S)^2$; and $x_S^{\mathsf T}Dx_S=\operatorname{vol}S/(\operatorname{vol}S)^2+\operatorname{vol}\bar S/(\operatorname{vol}\bar S)^2=1/\operatorname{vol}S+1/\operatorname{vol}\bar S$. The quotient is
$\operatorname{Ncut}(S)$. For the variational formula, substitute $y=D^{1/2}x$: the quotient becomes $y^{\mathsf T}\mathcal Ly/y^{\mathsf T}y$ with the \emph{normalised Laplacian}
$\mathcal L=D^{-1/2}LD^{-1/2}$, symmetric positive semidefinite, whose eigenvalues are those of $Lx=\lambda Dx$; its smallest eigenvalue is $0$ with eigenvector $D^{1/2}\mathbf1$, and the constraint
is $y\perp D^{1/2}\mathbf1$. By Courant--Fischer the minimum over this hyperplane is the second eigenvalue. The vectors $x_S$ are admissible, which gives the inequality.
\end{proof}

So $\lambda_2$ is a lower bound for the best normalised cut, computable in polynomial time, and the minimising eigenvector, a ``relaxed'' cut
taking real values, suggests a true cut: sort the vertices by its values and take the best of the $n-1$ threshold cuts, the \emph{sweep}. That the sweep is
also nearly optimal, $\operatorname{Ncut}\le2\sqrt{2\lambda_2}$ for the best threshold cut, is a Cheeger inequality, proved in Tome~C1, Theorem~9.9 for regular graphs; the normalised version is
proved the same way (Chung, 1997). In practice one computes the eigenvector by Lanczos on $\mathcal L$ (or on $2I-\mathcal L$, to make it the top one), with a cost
proportional to the number of edges. With $c$ clusters one takes the $c$ smallest eigenvectors as coordinates and runs $k$-means on them (Ng, Jordan and Weiss, 2002).
Figure~\ref{fig:moons} shows the method separating two interleaved half-moons, where any straight line fails.

\begin{figure}[ht]
\centering
\fig{ch02-moons}
\caption{Two interleaved half-moons of $300$ noisy points, coloured by the sign of the second eigenvector of the normalised Laplacian of their $10$-nearest-neighbour
graph. The sweep cut recovers the two moons without a single error, and its normalised cut is only $1.23$ times the lower bound $\lambda_2\approx1.0\times10^{-4}$.}
\label{fig:moons}
\end{figure}

\section{Nets and Gaussian matrices}\label{sec:nets}

We now turn to the second way of making a problem small. The Johnson--Lindenstrauss lemma (Tome~10, Theorem~8.10) says that a $k\times n$ matrix $S$ with independent
$N(0,\frac1k)$ entries preserves the squared lengths of $N$ given vectors up to a factor $1\pm\varepsilon$ as soon as $k\gtrsim\varepsilon^{-2}\log N$. For least squares we need
all the vectors of a $d$-dimensional subspace at once, an infinite set. The way from finite to infinite is the $\varepsilon$-net.

\begin{definition}\label{def:net}
A subset $\mathcal N$ of the unit sphere $\mathbb S^{d-1}\subset\R^d$ is an \emph{$\varepsilon$-net} if every point of the sphere is within distance $\varepsilon$ of some point of $\mathcal N$.
\end{definition}

\begin{lemma}[size of nets]\label{lem:netsize}
For $0<\varepsilon\le1$, $\mathbb S^{d-1}$ has an $\varepsilon$-net with at most $(1+2/\varepsilon)^d$ points.
\end{lemma}
\begin{proof}
Let $\mathcal N$ be a maximal subset of the sphere whose points are at mutual distance greater than $\varepsilon$ (it is finite by compactness, as the argument shows). By maximality
it is an $\varepsilon$-net. The open balls of radius $\varepsilon/2$ around its points are disjoint and contained in the ball of radius $1+\varepsilon/2$ around $0$. Comparing
Lebesgue measures, which scale like the $d$th power of the radius (Tome~6, Chapter~8), $|\mathcal N|(\varepsilon/2)^d\le(1+\varepsilon/2)^d$.
\end{proof}

\begin{lemma}[nets control norms]\label{lem:netnorm}
Let $B$ be a symmetric $d\times d$ matrix and $\mathcal N$ an $\varepsilon$-net of $\mathbb S^{d-1}$ with $\varepsilon<\frac12$. Then
$\|B\|=\max_{x\in\mathbb S^{d-1}}|x^{\mathsf T}Bx|\le(1-2\varepsilon)^{-1}\max_{y\in\mathcal N}|y^{\mathsf T}By|$.
\end{lemma}
\begin{proof}
The equality is the spectral theorem. Take a unit $x$ with $|x^{\mathsf T}Bx|=\|B\|$ and $y\in\mathcal N$ with $\|x-y\|\le\varepsilon$. Then
$|x^{\mathsf T}Bx-y^{\mathsf T}By|=|(x-y)^{\mathsf T}Bx+y^{\mathsf T}B(x-y)|\le2\varepsilon\|B\|$, so $\|B\|\le|y^{\mathsf T}By|+2\varepsilon\|B\|$.
\end{proof}

\begin{definition}\label{def:embedding}
A $k\times n$ matrix $S$ is a \emph{subspace embedding} with distortion $\varepsilon$ for a subspace $E\subseteq\R^n$ if $(1-\varepsilon)\|v\|^2\le\|Sv\|^2\le(1+\varepsilon)\|v\|^2$ for all $v\in E$.
\end{definition}

\begin{theorem}[Gaussian subspace embeddings]\label{thm:embedding}
Let $E\subseteq\R^n$ be a subspace of dimension $d$, $0<\varepsilon,\delta<1$, and $S$ a $k\times n$ matrix with independent $N(0,\frac1k)$ entries, where
\[
k\ge\frac{24}{\varepsilon^2}\Bigl(d\ln9+\ln\frac2\delta\Bigr).
\]
Then $S$ is a subspace embedding for $E$ with distortion $\varepsilon$ with probability at least $1-\delta$. Equivalently, if $G$ is a $k\times d$ matrix with independent
$N(0,\frac1k)$ entries, all singular values of $G$ lie in $[\sqrt{1-\varepsilon},\sqrt{1+\varepsilon}]$ with probability at least $1-\delta$.
\end{theorem}
\begin{proof}
Let $U$ be an $n\times d$ matrix whose columns are an orthonormal basis of $E$. Each row $s_i$ of $\sqrt kS$ is a standard Gaussian vector of $\R^n$, so $U^{\mathsf T}s_i$ is Gaussian
with covariance $U^{\mathsf T}U=I_d$ (Tome~3, Theorem~6.15), and the rows are independent: $G=SU$ has independent $N(0,\frac1k)$ entries. For $v=Ux$, $\|v\|=\|x\|$ and
$\|Sv\|=\|Gx\|$, so both statements say $\|G^{\mathsf T}G-I\|\le\varepsilon$. For a fixed unit $y\in\R^d$, $k\|Gy\|^2$ is a sum of $k$ squares of independent standard Gaussians
(as in Tome~10, Theorem~8.10), and Tome~10, Lemma~8.9 gives, for $0<t<1$, using $\frac{t^2}2-\frac{t^3}3\ge\frac{t^2}6$,
\[
P\bigl(|\|Gy\|^2-1|\ge t\bigr)\le e^{-kt^2/6}+e^{-kt^2/4}\le2e^{-kt^2/6}.
\]
Let $\mathcal N$ be a $\frac14$-net of $\mathbb S^{d-1}$ with at most $9^d$ points (Lemma~\ref{lem:netsize}). With $t=\varepsilon/2$, the union bound shows that
$\max_{y\in\mathcal N}|y^{\mathsf T}(G^{\mathsf T}G-I)y|<\varepsilon/2$ except with probability at most $2\cdot9^de^{-k\varepsilon^2/24}\le\delta$. On that event,
Lemma~\ref{lem:netnorm} with $\varepsilon=\frac14$ gives $\|G^{\mathsf T}G-I\|\le2\cdot\frac\varepsilon2=\varepsilon$.
\end{proof}

The number of rows depends on the dimension $d$ of the subspace, not on the ambient $n$: a subspace of dimension $50$ in $\R^{10^9}$ survives a projection to a few
thousand dimensions with all its geometry. The constant $24$ is generous. Gordon's theorem, proved with Gaussian comparison inequalities that we shall not need, says that the singular values of $G$ lie
within about $\sqrt{d/k}$ of $1$, so that $k\approx4d/\varepsilon^2$ rows already suffice for large $d$; the notebook measures it.

\section{Sketching least squares and low-rank approximation}\label{sec:sketching}

\emph{Least squares.} To minimise $\|Ax-b\|$ for a tall $n\times d$ matrix $A$, solve instead the small problem $\min_x\|S(Ax-b)\|$.

\begin{theorem}[sketch and solve; Sarl\'os, 2006]\label{thm:sketchsolve}
Let $S$ be a subspace embedding with distortion $\varepsilon$ for the span of the columns of $A$ and $b$, and $\tilde x$ a minimiser of $\|S(Ax-b)\|$. Then
\[
\|A\tilde x-b\|\le\sqrt{\frac{1+\varepsilon}{1-\varepsilon}}\ \min_x\|Ax-b\|.
\]
\end{theorem}
\begin{proof}
For every $x$, $Ax-b$ lies in the subspace. With $x_*$ a minimiser of $\|Ax-b\|$,
$(1-\varepsilon)\|A\tilde x-b\|^2\le\|S(A\tilde x-b)\|^2\le\|S(Ax_*-b)\|^2\le(1+\varepsilon)\|Ax_*-b\|^2$.
\end{proof}

By Theorem~\ref{thm:embedding}, $k=O((d+\log\frac1\delta)/\varepsilon^2)$ rows suffice. A dense Gaussian $S$ costs $O(knd)$ to apply, no better than solving the problem directly;
the gain comes from structured embeddings that are as good: the subsampled randomised Hadamard transform, applied by a fast transform in $O(nd\log n)$ (Ailon and
Chazelle, 2006; Sarl\'os, 2006), or the \emph{CountSketch}, with one random $\pm1$ per column, applied in time proportional to the number of nonzero entries of $A$
(Clarkson and Woodruff, 2013). The notebook compares them.

\emph{Low rank.} A matrix $A$ of size $m\times n$ whose singular values decay is well approximated by rank $r$, and the best approximation is given by the truncated
SVD (Eckart--Young, Tome~S1, Theorem~6.5): with $\sigma_1\ge\sigma_2\ge\dots$ the singular values,
\[
\min_{\operatorname{rank}B\le r}\|A-B\|_F^2=\tau_r^2:=\sum_{j>r}\sigma_j^2.
\]
Computing the SVD of a huge matrix is expensive. The \emph{randomised range finder} of Per-Gunnar Martinsson, Vladimir Rokhlin and Mark Tygert (2006), analysed by Nathan
Halko, Martinsson and Joel Tropp in 2011, does this: draw an $n\times k$ Gaussian matrix $\Omega$ with $k$ a bit larger than $r$, form $Y=A\Omega$ (one pass over $A$), compute an
orthonormal basis $Q$ of its columns, and approximate $A\approx Q(Q^{\mathsf T}A)$. The small matrix $Q^{\mathsf T}A$ has $k$ rows, and its SVD gives an approximate SVD of $A$.

\begin{think}
Why should $A\Omega$ capture the top singular directions? Write $\Omega$ in the basis of right singular vectors: what does $A$ do to the components of each column of
$\Omega$ along $v_1,\dots,v_r$, and along the others?
\end{think}

\begin{theorem}[randomised range finder]\label{thm:rangefinder}
Let $A$ be $m\times n$, $1\le r<k\le n$, $\Omega$ an $n\times k$ matrix with independent $N(0,1)$ entries, $Y=A\Omega$ and $P_Y$ the orthogonal projection onto the column space of $Y$.
\begin{enumerate}
\item (Deterministic bound.) Let $A=U\Sigma V^{\mathsf T}$ be a full SVD, $V=(V_1\ V_2)$ with $V_1$ of size $n\times r$, $\Omega_1=V_1^{\mathsf T}\Omega$, $\Omega_2=V_2^{\mathsf T}\Omega$, and $\Sigma_2$
the part of $\Sigma$ carrying $\sigma_{r+1},\sigma_{r+2},\dots$. If $\Omega_1$ has rank $r$, then $\|(I-P_Y)A\|_F^2\le\tau_r^2+\|\Sigma_2\Omega_2\Omega_1^+\|_F^2$, where $\Omega_1^+$ is the
pseudo-inverse (Tome~S1, Theorem~6.7).
\item (Average.) $\Omega_1$ and $\Omega_2$ are independent with independent $N(0,1)$ entries, and if $k\ge96(r\ln9+\ln\frac2\delta)$, then with probability at least $1-\delta$ over
$\Omega_1$, the conditional expectation over $\Omega_2$ satisfies $\mathbb E\bigl[\|(I-P_Y)A\|_F^2\,\big|\,\Omega_1\bigr]\le\bigl(1+\frac{2r}k\bigr)\tau_r^2$.
\end{enumerate}
\end{theorem}
\begin{proof}
(1) The columns of $P_YA$ are the closest points to those of $A$ in the column space of $Y$, so $\|(I-P_Y)A\|_F\le\|A-X\|_F$ for every $X$ whose columns lie in that space. Choose
$X=Y\Omega_1^+V_1^{\mathsf T}$. Writing $U=(U_1\ U_2)$ and $\Sigma=\begin{psmallmatrix}\Sigma_1&0\\0&\Sigma_2\end{psmallmatrix}$ conformally,
$Y=U\Sigma V^{\mathsf T}\Omega=U_1\Sigma_1\Omega_1+U_2\Sigma_2\Omega_2$, and since $\Omega_1\Omega_1^+=I_r$ (full row rank),
\[
X=U_1\Sigma_1V_1^{\mathsf T}+U_2\Sigma_2\Omega_2\Omega_1^+V_1^{\mathsf T},\qquad A-X=U_2\Sigma_2V_2^{\mathsf T}-U_2\Sigma_2\Omega_2\Omega_1^+V_1^{\mathsf T},
\]
using $A=U_1\Sigma_1V_1^{\mathsf T}+U_2\Sigma_2V_2^{\mathsf T}$. As $U_2$ has orthonormal columns and the two terms are orthogonal in the Frobenius inner product
($\operatorname{tr}(V_2\Sigma_2^{\mathsf T}\Sigma_2\Omega_2\Omega_1^+V_1^{\mathsf T})=\operatorname{tr}(V_1^{\mathsf T}V_2\cdots)=0$ since $V_1^{\mathsf T}V_2=0$), we get
$\|A-X\|_F^2=\|\Sigma_2\|_F^2+\|\Sigma_2\Omega_2\Omega_1^+\|_F^2$, and $\|\Sigma_2\|_F^2=\tau_r^2$.

(2) The columns of $V^{\mathsf T}\Omega$ are independent standard Gaussian vectors, since $V^{\mathsf T}$ is orthogonal; so its entries are independent $N(0,1)$, and $\Omega_1$, $\Omega_2$
are disjoint blocks of it. For a fixed matrix $M$ and $G$ with independent $N(0,1)$ entries, the entry $(i,j)$ of $\Sigma_2GM$ is $\sigma_{r+i}\sum_lg_{il}M_{lj}$, with mean $0$ and variance
$\sigma_{r+i}^2\sum_lM_{lj}^2$; summing, $\mathbb E\|\Sigma_2GM\|_F^2=\|\Sigma_2\|_F^2\|M\|_F^2$. Apply this with $G=\Omega_2$ and $M=\Omega_1^+$, independent of $\Omega_2$. By
Theorem~\ref{thm:embedding} with $\varepsilon=\frac12$, applied to the $k\times r$ matrix $\Omega_1^{\mathsf T}/\sqrt k$, with probability at least $1-\delta$ all singular values of $\Omega_1$ are
at least $\sqrt{k/2}$; then $\Omega_1$ has rank $r$ and $\|\Omega_1^+\|_F^2=\sum_{i\le r}\sigma_i(\Omega_1)^{-2}\le2r/k$. Combine with (1).
\end{proof}

Read the bound: the randomised method loses a factor $1+\frac{2r}k$ against the best possible, with a constant in the condition on $k$ that is far too pessimistic. Halko, Martinsson and
Tropp proved, by computing $\mathbb E\|\Omega_1^+\|_F^2=\frac r{k-r-1}$ exactly (a property of inverse Wishart matrices that we do not prove), that $k=r+p$ with $p\ge2$ already gives
$\mathbb E\|(I-P_Y)A\|_F^2\le(1+\frac r{p-1})\tau_r^2$; ten extra columns are plenty in practice. When the singular values decay slowly, one replaces $Y=A\Omega$ by $(AA^{\mathsf T})^qA\Omega$,
which has the same singular vectors and singular values $\sigma_j^{2q+1}$, much more decaying. Figure~\ref{fig:rsvd} shows the error on a test image.

\begin{figure}[ht]
\centering
\fig{ch02-rsvd}
\caption{Randomised range finder on a $512\times512$ test image (the notebook's synthetic picture): relative Frobenius error $\|(I-P_Y)A\|_F/\|A\|_F$ with $k=r+10$ columns, without and with one
power iteration ($q=1$), against the best rank-$r$ error $\tau_r/\|A\|_F$ of the truncated SVD. Without power iteration the loss is a factor $1.1$ to $1.5$; with it,
the $r+10$ columns even beat the best rank-$r$ approximation.}
\label{fig:rsvd}
\end{figure}

\begin{beauty}
Conjugate gradients and the randomised range finder look like opposites, one deterministic and adaptive, the other random and oblivious. Both rest on the same
geometric fact: the information needed lives in a small subspace, and once we are confined to it, the dimension of the ambient space disappears from the
bounds. Krylov finds the subspace by the optimal polynomial; the Gaussian matrix finds it by the concentration of measure that makes a random $k$-dimensional
view of a $d$-dimensional space almost an isometry.
\end{beauty}

\section{Exercises}

\begin{exercise}[$\star$]\label{ex:cgcount}
Using Theorem~\ref{thm:cgconv}(2), estimate the number of conjugate gradient steps needed to reduce the energy error by $10^{-6}$ when $\kappa=10^4$, and compare with the
$\frac{\kappa}2\ln(10^6)$ steps of an iteration contracting by $\frac{\kappa-1}{\kappa+1}$ per step.
\end{exercise}

\begin{exercise}[$\star$]\label{ex:krylovdim}
Show that $\dim\mathcal K_k(A,b)=\min(k,m)$, where $m$ is the degree of the minimal polynomial of $b$ with respect to $A$ (the monic polynomial $p$ of least degree with
$p(A)b=0$). Deduce that conjugate gradients stops with $x_m=x_*$, and relate $m$ to the eigenvectors needed to express $b$.
\end{exercise}

\begin{exercise}[$\star\star$]\label{ex:cgnormal}
Let $A$ be invertible but not symmetric. Show that conjugate gradients applied to $A^{\mathsf T}Ax=A^{\mathsf T}b$ (CGNR) minimises $\|b-Ax\|$ over $x\in\mathcal K_k(A^{\mathsf T}A,A^{\mathsf T}b)$,
and that its rate is governed by $\kappa(A)$, not $\sqrt{\kappa(A)}$.
\end{exercise}

\begin{exercise}[$\star\star$]\label{ex:steepest}
Steepest descent for $\frac12x^{\mathsf T}Ax-b^{\mathsf T}x$ moves along $r_k$ with the exact minimising step. Show that it is the conjugate gradient method restarted after every step, and
that for $A=\operatorname{diag}(1,\kappa)$ and a suitable starting point the energy error is multiplied by exactly $\frac{\kappa-1}{\kappa+1}$ at every step.
\end{exercise}

\begin{exercise}[$\star\star$]\label{ex:chebopt}
Let $0<\mu<L$ and $k\ge1$. Show that every $p\in\mathcal P_k$ with $p(0)=1$ satisfies $\max_{[\mu,L]}|p|\ge1/T_k(x_0)$, with $x_0=\frac{L+\mu}{L-\mu}$, so that the
polynomial $p_k$ of~\eqref{eq:chebpoly} is optimal. So Theorem~\ref{thm:cgconv}(2) is the best bound that depends on $\kappa$ alone.
\end{exercise}

\begin{exercise}[$\star$]\label{ex:netlower}
Let $0<\varepsilon\le1$. Show that every $\varepsilon$-net of $\mathbb S^{d-1}$ has at least $2^{-d}\varepsilon^{-(d-1)}$ points. For $d=2$, show that the minimal size of an
$\varepsilon$-net is $\lceil\pi/(2\arcsin(\varepsilon/2))\rceil$.
\end{exercise}

\begin{exercise}[$\star$]\label{ex:normlower}
Let $G$ be $k\times d$ with independent $N(0,\frac1k)$ entries. Show that $\mathbb E\|G\|^2\ge1+\frac{d-1}k$. Deduce that a Gaussian sketch whose largest singular value has
mean square distortion $\mathbb E\|G\|^2-1\le\varepsilon$ needs $k\ge(d-1)/\varepsilon$ rows: the dimension $d$ in Theorem~\ref{thm:embedding} cannot be avoided.
\end{exercise}

\begin{exercise}[$\star\star$]\label{ex:sketchsol}
In Theorem~\ref{thm:sketchsolve}, show that the sketched solution satisfies $\|A(\tilde x-x_*)\|^2\le\frac{2\varepsilon}{1-\varepsilon}\|Ax_*-b\|^2$. Can $\|\tilde x-x_*\|$ itself be large? Give an example.
\end{exercise}

\begin{exercise}[$\star$]\label{ex:exactlowrank}
Show that if $A$ has rank $r$ and $k\ge r$, the randomised range finder recovers $A$ exactly, $P_YA=A$, with probability $1$.
\end{exercise}

\begin{exercise}[$\star\star$]\label{ex:ncutwalk}
Let $X_0,X_1$ be two steps of the random walk on the weighted graph started from its stationary distribution $\pi_i=d_i/\operatorname{vol}V$. Show that
$\operatorname{Ncut}(S)=P(X_1\notin S\mid X_0\in S)+P(X_1\in S\mid X_0\notin S)$ (Meil\u{a} and Shi, 2001).
\end{exercise}

\begin{exercise}[$\star\star$]\label{ex:ncutpath}
For the path with $n$ vertices and unit weights, compute $\lambda_2$ of the normalised Laplacian and the best normalised cut, and check Theorem~\ref{thm:ncut}.
\end{exercise}

\begin{exercise}[$\star\star$]\label{ex:ghosts}
Run the Lanczos recurrence in floating point on $A=\operatorname{diag}(1,2,\dots,100)$ with $b=\mathbf1$ for $200$ steps, and explain why some eigenvalues of $T_{200}$ appear twice.
Show that if the $q_j$ were exactly orthonormal, $T_k$ could not have more than $100$ distinct eigenvalues, nor any eigenvalue repeated.
\end{exercise}

\section{Hints}

\begin{hint}{ex:cgcount}
Solve $2(1-2/\sqrt\kappa)^k\le10^{-6}$ approximately.
\end{hint}
\begin{hint}{ex:krylovdim}
$A^kb\in\mathcal K_k$ exactly when some polynomial of degree $k$ annihilates $b$.
\end{hint}
\begin{hint}{ex:cgnormal}
$\|x-x_*\|_{A^{\mathsf T}A}=\|Ax-b\|$, and $\kappa(A^{\mathsf T}A)=\kappa(A)^2$.
\end{hint}
\begin{hint}{ex:steepest}
Start from $x_0$ with $A$-energy components such that the iterates alternate between two directions; try $x_*-x_0=(\kappa,1)$.
\end{hint}
\begin{hint}{ex:chebopt}
$T_k$ equioscillates: it takes the values $\pm1$ alternately at $k+1$ points of $[-1,1]$. If $q$ did strictly better, $p_k-q$ would have $k$ roots in $(\mu,L)$ and a root at $0$.
\end{hint}
\begin{hint}{ex:netlower}
The caps of radius $\varepsilon$ around the net points cover the sphere. To bound the area of a cap, thicken it into a shell of width $\varepsilon$ and put the shell in a ball of radius $2\varepsilon$.
\end{hint}
\begin{hint}{ex:normlower}
Let $g_1$ be the first column of $G$ and $u=g_1/\|g_1\|$. Then $\|G\|^2\ge\|G^{\mathsf T}u\|^2$, and $u$ is independent of the other columns.
\end{hint}
\begin{hint}{ex:sketchsol}
$Ax_*-b$ is orthogonal to the column space of $A$, so $\|A\tilde x-b\|^2=\|A(\tilde x-x_*)\|^2+\|Ax_*-b\|^2$.
\end{hint}
\begin{hint}{ex:exactlowrank}
In Theorem~\ref{thm:rangefinder}(1), $\Sigma_2=0$; you only need $\Omega_1$ to have rank $r$.
\end{hint}
\begin{hint}{ex:ncutwalk}
$P(X_0=i,X_1=j)=w_{ij}/\operatorname{vol}V$.
\end{hint}
\begin{hint}{ex:ncutpath}
The eigenvectors of the normalised Laplacian of a path are cosines; the best cut is in the middle.
\end{hint}
\begin{hint}{ex:ghosts}
A Ritz value that has converged is found again by the Krylov process once rounding errors reintroduce a component along its eigenvector.
\end{hint}

\section{Solutions}

\begin{solution}{ex:cgcount}
\idea Take logarithms in the bound.
\writeup We need $k\ln\frac{\sqrt\kappa+1}{\sqrt\kappa-1}\ge\ln(2\times10^6)$. With $\sqrt\kappa=100$, $\ln\frac{101}{99}\approx0.0200$, so $k\ge14.5/0.0200\approx726$. The other iteration needs
$k\ln\frac{\kappa+1}{\kappa-1}\ge\ln10^6$, with $\ln\frac{10001}{9999}\approx2\times10^{-4}$, so $k\approx69\,000$: a hundred times more, the factor $\sqrt\kappa$.
\end{solution}

\begin{solution}{ex:krylovdim}
\idea The Krylov vectors stop being independent exactly when a polynomial relation appears.
\writeup If $b,Ab,\dots,A^{j-1}b$ are independent but $A^jb$ depends on them, then a monic polynomial of degree $j$ annihilates $b$, and conversely; so the dimensions grow by one until
$k=m$, where $A^mb\in\mathcal K_m$, and then $A\mathcal K_m\subseteq\mathcal K_m$ so all later spaces equal $\mathcal K_m$. As $\mathcal K_m$ is invariant and $A$ invertible, $A$ maps it onto itself
and $x_*=A^{-1}b\in\mathcal K_m$, so $x_m=x_*$ by Definition~\ref{def:cg}. If $b=\sum c_iv_i$ in an eigenbasis, $p(A)b=\sum c_ip(\lambda_i)v_i$ vanishes iff $p$ vanishes at the distinct
eigenvalues $\lambda_i$ with $c_i\ne0$: $m$ is the number of distinct eigenvalues ``seen'' by $b$.
\end{solution}

\begin{solution}{ex:cgnormal}
\idea Rewrite the energy norm of the normal equations.
\writeup $A^{\mathsf T}A$ is symmetric positive definite with solution $x_*$ of $Ax=b$, and $\|x-x_*\|^2_{A^{\mathsf T}A}=\|A(x-x_*)\|^2=\|Ax-b\|^2$. So the conjugate gradient iterate for the normal
equations minimises the residual norm over $\mathcal K_k(A^{\mathsf T}A,A^{\mathsf T}b)$. Its convergence is governed by $\kappa(A^{\mathsf T}A)=\kappa(A)^2$, so the rate
$\frac{\sqrt{\kappa(A^{\mathsf T}A)}-1}{\sqrt{\kappa(A^{\mathsf T}A)}+1}=\frac{\kappa(A)-1}{\kappa(A)+1}$ involves $\kappa(A)$ itself. (One never forms $A^{\mathsf T}A$: each step multiplies by $A$ and by
$A^{\mathsf T}$.)
\end{solution}

\begin{solution}{ex:steepest}
\idea The first step of conjugate gradients is a steepest descent step, so steepest descent is CG with memory one. For the rate, find a starting error that reproduces itself up to scale.
\writeup Started at $x_{k}$, the first CG step minimises the energy over $x_k+\operatorname{span}(r_k)$, which is the exact line search along the gradient direction: steepest descent restarts CG at each step.
Let $A=\operatorname{diag}(1,\kappa)$ and $e=x-x_*=(a,c)$. Then $r=-Ae=-(a,\kappa c)$, the step is $t=\frac{r^{\mathsf T}r}{r^{\mathsf T}Ar}=\frac{a^2+\kappa^2c^2}{a^2+\kappa^3c^2}$, and the new error is
$e'=e+tr=(a(1-t),c(1-\kappa t))$. With $e=(\kappa,1)$: $t=\frac{2\kappa^2}{\kappa^2+\kappa^3}=\frac2{1+\kappa}$, $e'=\frac{\kappa-1}{\kappa+1}(\kappa,-1)$. The next step from $(\kappa,-1)$ is symmetric and returns to
a multiple of $(\kappa,1)$. Each step multiplies $e$, hence $\|e\|_A$, by $\frac{\kappa-1}{\kappa+1}$: the zigzag of Figure~\ref{fig:cgconv} and of the cover of this book.
\end{solution}

\begin{solution}{ex:chebopt}
\idea Equioscillation plus root counting: a polynomial that is smaller everywhere would have to cross the Chebyshev one too many times.
\writeup Put $M=1/T_k(x_0)$. Since $T_k(\cos t)=\cos kt$, $T_k$ takes the value $(-1)^j$ at $\cos(j\pi/k)$, $j=0,\dots,k$. Through the affine map of~\eqref{eq:chebpoly}, $p_k$ takes the values
$(-1)^jM$ at $k+1$ points $\mu\le\xi_0<\xi_1<\dots<\xi_k\le L$ (with $\xi_0=\mu$, where $p_k=M>0$). Suppose $q\in\mathcal P_k$ has $q(0)=1$ and $\max_{[\mu,L]}|q|<M$. Then
$r=p_k-q$ satisfies $(-1)^jr(\xi_j)=M-(-1)^jq(\xi_j)>0$: $r$ changes sign on each of the $k$ intervals $(\xi_j,\xi_{j+1})$, so it has $k$ distinct roots in $(\mu,L)$, and also
$r(0)=1-1=0$ with $0<\mu$. A polynomial of degree at most $k$ with $k+1$ distinct roots is zero, so $q=p_k$, whose maximum is $M$: a contradiction. Hence $\max_{[\mu,L]}|q|\ge M$
for every admissible $q$. Since the eigenvalues of $A$ can fill $[\mu,L]$ as densely as we like, no bound of the form $\max_{[\mu,L]}|p|$ beats Theorem~\ref{thm:cgconv}(2).
\end{solution}

\begin{solution}{ex:netlower}
\idea Count by area: each net point covers only a small cap, and a cap is small because a thin shell over it fits in a small ball.
\writeup Let $C(y)=\{x\in\mathbb S^{d-1}:\|x-y\|\le\varepsilon\}$ and let $\sigma$ be surface measure, so that $\sigma(\mathbb S^{d-1})=d\,\omega_d$ with $\omega_d$ the volume of the unit ball. The caps $C(y)$,
$y\in\mathcal N$, cover the sphere, so $|\mathcal N|\ge\sigma(\mathbb S^{d-1})/\max_y\sigma(C(y))$. The shell $\{tx:1-\varepsilon\le t\le1,\ x\in C(y)\}$ has volume
$\sigma(C(y))\,(1-(1-\varepsilon)^d)/d\ge\sigma(C(y))\,\varepsilon/d$ (polar coordinates), and each of its points satisfies $\|tx-y\|\le\|tx-x\|+\|x-y\|\le2\varepsilon$, so it lies in a ball of volume
$\omega_d(2\varepsilon)^d$. Hence $\sigma(C(y))\le d\,\omega_d2^d\varepsilon^{d-1}$ and $|\mathcal N|\ge2^{-d}\varepsilon^{-(d-1)}$. For fixed $\varepsilon<\frac12$ this grows exponentially in $d$, as
Lemma~\ref{lem:netsize} does: the union bound over a net really pays a price of order $d$ in the exponent.

For $d=2$, $\|x-y\|\le\varepsilon$ on the unit circle means that the angle between $x$ and $y$ is at most $\alpha=2\arcsin(\varepsilon/2)$, so each cap is an arc of length $2\alpha$, and $N$ arcs cover
the circle only if $2N\alpha\ge2\pi$, that is $N\ge\pi/\alpha$. Conversely, $N$ equally spaced points leave every point within angle $\pi/N$ of one of them, and $\pi/N\le\alpha$ as soon as
$N\ge\pi/\alpha$. The minimal size is $\lceil\pi/(2\arcsin(\varepsilon/2))\rceil$.
\end{solution}

\begin{solution}{ex:normlower}
\idea Test $G^{\mathsf T}$ on the direction of its own first column, which sees the full length of that column and, in addition, a random projection of every other column.
\writeup Let $g_1,\dots,g_d\in\R^k$ be the columns and $u=g_1/\|g_1\|$ (defined with probability $1$). Then $\|G\|^2=\|G^{\mathsf T}\|^2\ge\|G^{\mathsf T}u\|^2=\|g_1\|^2+\sum_{j\ge2}(g_j^{\mathsf T}u)^2$.
Now $\mathbb E\|g_1\|^2=k\cdot\frac1k=1$, and for $j\ge2$, $g_j$ is independent of $u$; conditionally on $u$, $g_j^{\mathsf T}u$ is $N(0,\frac1k)$ because $u$ is a unit vector and the law of $g_j$ is
rotation invariant. So $\mathbb E(g_j^{\mathsf T}u)^2=\frac1k$ and $\mathbb E\|G\|^2\ge1+\frac{d-1}k$. If $\mathbb E\|G\|^2-1\le\varepsilon$, then $\frac{d-1}k\le\varepsilon$, that is
$k\ge(d-1)/\varepsilon$. (Gordon's theorem improves this to the true order $k\gtrsim d/\varepsilon^2$, which the notebook observes.)
\end{solution}

\begin{solution}{ex:sketchsol}
\idea Pythagoras separates the error of $\tilde x$ into a part we control and the unavoidable residual.
\writeup Since $Ax_*-b\perp\operatorname{range}A$, $\|A\tilde x-b\|^2=\|A(\tilde x-x_*)\|^2+\|Ax_*-b\|^2$. By the proof of Theorem~\ref{thm:sketchsolve} the left side is at most $\frac{1+\varepsilon}{1-\varepsilon}\|Ax_*-b\|^2$, so
$\|A(\tilde x-x_*)\|^2\le\frac{2\varepsilon}{1-\varepsilon}\|Ax_*-b\|^2$. The vector $\tilde x-x_*$ itself can be large when $A$ is ill-conditioned: $\|\tilde x-x_*\|\le\|A(\tilde x-x_*)\|/\sigma_{\min}(A)$, and for
$A=\operatorname{diag}(1,\eta)$ padded with zero rows and $b$ with a residual of size $1$, the second coordinate of $\tilde x$ can be off by about $\sqrt{2\varepsilon}/\eta$. Also, if $b$ is in the range of $A$,
$\tilde x=x_*$ exactly.
\end{solution}

\begin{solution}{ex:exactlowrank}
\idea Use the deterministic bound with $\Sigma_2=0$.
\writeup If $\operatorname{rank}A=r$, then $\sigma_{r+1}=\dots=0$, so $\Sigma_2=0$ and $\tau_r=0$. By Theorem~\ref{thm:rangefinder}(1), $\|(I-P_Y)A\|_F=0$ as soon as $\Omega_1=V_1^{\mathsf T}\Omega$ ($r\times k$ Gaussian)
has rank $r$. The set of $r\times k$ matrices of rank $<r$ is the zero set of the $r\times r$ minors, a proper algebraic subset, of Lebesgue measure zero; the Gaussian law has a density, so it
has probability $0$.
\end{solution}

\begin{solution}{ex:ncutwalk}
\idea Stationary edge probabilities are proportional to the weights.
\writeup With $\pi_i=d_i/\operatorname{vol}V$ and transition probabilities $w_{ij}/d_i$, $P(X_0=i,X_1=j)=w_{ij}/\operatorname{vol}V$. So $P(X_0\in S,X_1\notin S)=\operatorname{cut}(S,\bar S)/\operatorname{vol}V$
and $P(X_0\in S)=\operatorname{vol}S/\operatorname{vol}V$, hence $P(X_1\notin S\mid X_0\in S)=\operatorname{cut}(S,\bar S)/\operatorname{vol}S$; similarly for $\bar S$, and the sum is $\operatorname{Ncut}(S)$. A good
normalised cut is a set from which the random walk rarely escapes.
\end{solution}

\begin{solution}{ex:ncutpath}
\idea Guess cosine eigenvectors for the normalised problem $Lx=\lambda Dx$; for the cut, a single edge is cut.
\writeup On the path with vertices $1,\dots,n$, the two ends have degree $1$ and the other vertices degree $2$, and $Lx=\lambda Dx$ reads $2x_i-x_{i-1}-x_{i+1}=2\lambda x_i$ inside and $x_1-x_2=\lambda x_1$ at the end. The vectors $x_i=\cos(\pi j(i-1)/(n-1))$ satisfy
these with $\lambda=1-\cos(\pi j/(n-1))$: inside, $x_{i-1}+x_{i+1}=2\cos(\pi j/(n-1))x_i$, and at the ends one checks directly. So $\lambda_2=1-\cos\frac\pi{n-1}\approx\frac{\pi^2}{2(n-1)^2}$. A cut of a path removes at least one
edge; cutting the edge $(s,s+1)$ gives $\operatorname{vol}S=2s-1$ and $\operatorname{vol}\bar S=2(n-s)-1$, so
\[
\operatorname{Ncut}=\frac1{2s-1}+\frac1{2n-2s-1},
\]
minimal in the middle, about $\frac2n$. Indeed
$\lambda_2\approx\frac{4.93}{n^2}\le\frac2n$, and the relaxation loses a factor of order $n$ on the path. That is the worst case allowed by Cheeger's inequality, which bounds the best cut by a
constant times $\sqrt{\lambda_2}\approx2.2/n$, the order of $\frac2n$ (Chung, Chapter~2).
\end{solution}

\begin{solution}{ex:ghosts}
\idea In exact arithmetic the Krylov space has dimension at most $100$; in floating point the lost orthogonality lets converged eigenvectors re-enter.
\writeup With exact orthonormal $q_j$, $\mathcal K_k(A,b)$ has dimension at most $100$ (Exercise~\ref{ex:krylovdim}; here $b$ sees all $100$ eigenvalues), so the process breaks down with $\beta_{100}=0$ and $T_{100}$
has exactly the eigenvalues $1,\dots,100$; and a symmetric tridiagonal matrix with nonzero off-diagonal entries has distinct eigenvalues (an eigenvector is determined by its first entry through the
recurrence). In floating point, as soon as a Ritz value $\theta$ converges to an eigenvalue $\lambda$, the new vectors $q_j$ acquire, by rounding, components along the corresponding eigenvector $v$ that
are no longer removed (Paige's theorem relates the loss of orthogonality to convergence), and the process ``rediscovers'' $\lambda$: $T_k$ gets a second Ritz value near $\lambda$, a ghost. The notebook
shows the extreme eigenvalues $1$ and $100$ appearing several times after $200$ steps, while full reorthogonalisation removes the ghosts.
\end{solution}

\section*{Sources and further reading}

\begin{itemize}
\item \emph{A research paper to read.} N.~Halko, P.-G.~Martinsson and J.~A.~Tropp, \emph{Finding structure with randomness: probabilistic algorithms for constructing approximate matrix decompositions},
SIAM Review 53 (2011), 217--288. Read Sections~1 and~4 for the algorithms (our Section~\ref{sec:sketching}), then Section~9, whose Theorem~9.1 is our deterministic bound with the operator norm,
and Section~10, where the expectation of $\|\Omega_1^+\|_F^2$ is computed. Then compare with D.~P.~Woodruff, \emph{Sketching as a tool for numerical linear algebra}, Foundations and Trends in Theoretical
Computer Science 10 (2014), the reference for subspace embeddings and CountSketch.
\item \emph{The stories.} M.~R.~Hestenes and E.~Stiefel, \emph{Methods of conjugate gradients for solving linear systems}, Journal of Research of the National Bureau of Standards 49 (1952), 409--436 (the
epigraph is from its introduction, which also mentions the August 1951 symposium). D.~P.~O'Leary, \emph{Methods of conjugate gradients for solving linear systems}, in \emph{A Century of Excellence in
Measurements, Standards, and Technology}, NIST Special Publication 958 (2001), 81--85, for the Institute for Numerical Analysis, the eclipse of the 1960s and Reid's revival; G.~H.~Golub and
D.~P.~O'Leary, \emph{Some history of the conjugate gradient and Lanczos algorithms: 1948--1976}, SIAM Review 31 (1989), 50--102; J.~K.~Reid, \emph{On the method of conjugate gradients for the solution of
large sparse systems of linear equations}, in \emph{Large Sparse Sets of Linear Equations} (J.~K.~Reid, ed.), Academic Press, 1971, 231--254; J.~Dongarra and F.~Sullivan, \emph{Guest editors' introduction:
the top 10 algorithms}, Computing in Science and Engineering 2 (2000), 22--23.
\item C.~Lanczos, \emph{An iteration method for the solution of the eigenvalue problem of linear differential and integral operators}, Journal of Research of the National Bureau of Standards 45 (1950),
255--282; S.~Kaniel, \emph{Estimates for some computational techniques in linear algebra}, Mathematics of Computation 20 (1966), 369--378; C.~C.~Paige, \emph{The computation of eigenvalues and eigenvectors of
very large sparse matrices}, PhD thesis, University of London, 1971; Y.~Saad, \emph{On the rates of convergence of the Lanczos and the block-Lanczos methods}, SIAM Journal on Numerical Analysis 17 (1980),
687--706; A.~Greenbaum, \emph{Behavior of slightly perturbed Lanczos and conjugate-gradient recurrences}, Linear Algebra and its Applications 113 (1989), 7--63; D.~A.~Spielman and S.-H.~Teng,
\emph{Nearly-linear time algorithms for graph partitioning, graph sparsification, and solving linear systems}, STOC 2004, 81--90.
\item J.~Shi and J.~Malik, \emph{Normalized cuts and image segmentation}, IEEE Transactions on Pattern Analysis and Machine Intelligence 22 (2000), 888--905; M.~Meil\u{a} and J.~Shi, \emph{A random walks
view of spectral segmentation}, AISTATS 2001; A.~Y.~Ng, M.~I.~Jordan and Y.~Weiss, \emph{On spectral clustering: analysis and an algorithm}, NIPS 14 (2002), 849--856; F.~R.~K.~Chung, \emph{Spectral Graph
Theory}, AMS, 1997; U.~von~Luxburg, \emph{A tutorial on spectral clustering}, Statistics and Computing 17 (2007), 395--416.
\item T.~Sarl\'os, \emph{Improved approximation algorithms for large matrices via random projections}, FOCS 2006, 143--152; N.~Ailon and B.~Chazelle, \emph{Approximate nearest neighbors and the fast
Johnson--Lindenstrauss transform}, STOC 2006, 557--563; K.~L.~Clarkson and D.~P.~Woodruff, \emph{Low rank approximation and regression in input sparsity time}, STOC 2013, 81--90; P.-G.~Martinsson,
V.~Rokhlin and M.~Tygert, \emph{A randomized algorithm for the approximation of matrices}, Yale technical report 1361 (2006), and Applied and Computational Harmonic Analysis 30 (2011), 47--68.
R.~Vershynin, \emph{High-Dimensional Probability}, Cambridge University Press, 2018, Chapter~4 for nets and Gaussian matrices, Section~7.3 for Gordon's theorem.
\item For Krylov methods in general: Y.~Saad, \emph{Iterative Methods for Sparse Linear Systems}, second edition, SIAM, 2003; L.~N.~Trefethen and D.~Bau, \emph{Numerical Linear Algebra}, SIAM, 1997, Lectures 32 to 40.
\end{itemize}
